\section{Arithmetic language and constructive factorization}\label{sec:arith}

Throughout, $[P]$ denotes the bit of a decidable predicate. Every divisor
is positive. The value of $\mu$ at zero is immaterial and may be set to
zero by an outer test. Write $\mathbf1(n)=1$ and $\varepsilon(n)=[n=1]$ on
positive inputs. A function in a theorem schema means a fixed
primitive recursive (PR) definition with its numerical parameters free,
or a supplied finite table covering the computation.

\paragraph{The object theory.}
We use first-order intuitionistic arithmetic with defining equations
for PR functions, the successor axioms, and induction on quantifier-free
formulas. An integer can be represented by $(a,b)$ with value $a-b$.
Equality to $(c,d)$ is $a+d=b+c$. Multiplication uses
$(ac+bd,ad+bc)$. Rationals have a positive denominator and equality by
cross multiplication. Finite lists can be built by iterated pairing,
with length carried separately, so every access is a bounded recursion.
Polynomial coefficients, residue classes, tables of finite operations,
and tuples of rationals have codes of this kind.
For a decidable test $R(i,\mathbf u)$ and a PR bound $B$,
$\prod_{i<B}[R(i,\mathbf u)]=1$ expresses that every test passes.
The sum of the bits has positive sign if a witness exists, and the first
witness is found by a scan of $B$ entries. A proof by strong induction
on a decidable matrix uses ordinary induction on the product of all its
preceding bits. Finite sums below are reindexed by explicit maps or by
partitioning a bounded rectangle into its fibers.

The real notation abbreviates arithmetic statements about rational
approximants. For example, a tail estimate means a universal statement
in the precision and the two truncation indices, after clearing positive
denominators and adding the evaluation tolerance. Complex
inputs may be Gaussian rationals or fixed PR Cauchy names with a supplied
modulus, size bound and the indicated domain margins. Proofs about such
names are schemata in the defining algorithms. A finite input version
takes the needed samples and tests the finitely many error bounds in the
proof. No quantification over infinite sequences is part of iPRA.

\paragraph{Division and prime factors.}
For $n>0$, define $p\mid n$ by a zero remainder for $p>0$.
Primality of $p\ge2$ is the absence of divisors between $2$ and $p-1$.
Define $v_p(n)$ for $p\ge2$ by the largest $j\le n$ such that $p^j\mid n$.
The search is bounded because $2^j\ge j+1$. Every $n>1$ has a least
divisor $p>1$, obtained by scanning through $n$. A proper factor of
$p$ would contradict leastness, so $p$ is prime. Repeatedly divide the
current remainder by its least prime, stopping at $1$ and then padding
with idle steps. A nonidle step strictly decreases the remainder, hence
$n$ steps suffice. The product of extracted primes times the remainder
is always $n$.

For uniqueness, the Euclidean divisions of $a,b$ strictly decrease the
nonzero remainder. Back substitution through this bounded list gives
integers $u,v$ with $ua+vb=\gcd(a,b)$. Initially represent $a,b$ by
$(1,0),(0,1)$ and subtract quotient times the preceding representation
at each division. Thus $p\nmid a$ implies $ua+vp=1$ for prime $p$, so
$p\mid ab$ forces $p\mid b$. Apply this lemma successively along any
second list representing $n$, cancel the matching prime, and induct
on list length. This proves uniqueness up to permutation, additivity
$v_p(ab)=v_p(a)+v_p(b)$, and the bijection between divisors of $n$
and exponent vectors $0\le e_p\le v_p(n)$. Both directions of this
bijection are bounded computations (take valuations, or multiply prime
powers). These arguments can also start from ProveIt's valuation and
Bézout derivations.

\section{Cancellation, inversion, and counting}\label{sec:core}
Set
\[
 \omega(n)=\sum_{\substack{p\le n\\p\text{ prime}}}[v_p(n)>0],\quad
 \Omega(n)=\sum_{\substack{p\le n\\p\text{ prime}}}v_p(n),\quad
 S(n)=\prod_{\substack{p\le n\\p\text{ prime}}}[v_p(n)\le1].
\]
Define $\lambda(n)=(-1)^{\Omega(n)}$ and $\mu(n)=S(n)(-1)^{\omega(n)}$.
Thus $\mu(1)=1$, $\mu(p)=-1$, $\mu(p^a)=0$ for $a\ge2$.
A square divisor $t^2$ with $t>1$ supplies a prime $p\mid t$. Conversely
$p^2\mid n$ is already such a divisor. This proves
\[
 |\mu(n)|=\mu(n)^2=S(n),\qquad
 \mu(n)=[\omega(n)=\Omega(n)]\lambda(n).
\]
Indeed every positive exponent contributes one to $\omega$, and the
excess in $\Omega$ is a sum of nonnegative integers. If $a,b$ are
coprime their prime supports are disjoint, so both the squarefree
indicator and the sign factor multiply. If $\gcd(a,b)>1$, its least
prime divisor has square dividing $ab$, giving $\mu(ab)=0$.
All these alternatives are decidable. For instance $30=2\cdot3\cdot5$
is the least composite with $\mu=-1$. A squarefree composite with an odd
number of prime divisors needs at least three, and its three smallest
are at least $2,3,5$. The same argument gives $2310=2\cdot3\cdot5\cdot7\cdot11$
as the first integer with five distinct prime factors.

\begin{theorem}[Divisor cancellation]
For $n\ge1$,
\[
 \sum_{d\mid n}\mu(d)=\varepsilon(n),\qquad
 \sum_{d\mid n}|\mu(d)|=2^{\omega(n)}.
\]
If $n>1$, exactly $2^{\omega(n)-1}$ divisors have each nonzero sign.
\end{theorem}
\begin{proof}
List the distinct prime divisors of $n$ in increasing order. Its squarefree
divisors are in bijection with bit masks of length $\omega(n)$ by the
product of the selected primes. For $n>1$ flip the first bit. This is
a fixed-point-free involution reversing the parity of the number of
selected bits. Exactly one member of each pair has its first bit clear.
The other bits are unrestricted. For $n=1$ there is one empty mask.
Grouping the masks by their number of selected bits gives also
$\sum_j(-1)^j\binom rj=[r=0]$.
\end{proof}

Define the Dirichlet convolution
$(f*g)(n)=\sum_{d\mid n}f(d)g(n/d)$. It is commutative for scalar functions by
$d\leftrightarrow n/d$, associative because both bracketings enumerate
the same ordered triples with product $n$, and has unit $\varepsilon$.
For coprime $u,v$, the divisors of $uv$ are exactly $ab$ with $a\mid u$,
$b\mid v$. Factoring the two finite sums shows that convolution preserves
multiplicativity. The theorem says $\mu*\mathbf1=\varepsilon$.
It uniquely specifies $\mu$. Its value at $1$ must be $1$, and at
$n>1$ it is minus the sum over proper divisors. This triangular recursion
also computes $\mu$ without factoring, by building a table through $n$.
\begin{theorem}[Two-sided inversion]
For fixed integer-valued PR functions, $F(n)=\sum_{d\mid n}f(d)$ for
all $n\ge1$ if and only if
$f(n)=\sum_{d\mid n}\mu(d)F(n/d)$ for all $n\ge1$.
The same holds up to any fixed bound and in a coded abelian group.
\end{theorem}
\begin{proof}
Convolve the forward equality with $\mu$, or the reverse one with
$\mathbf1$, and use associativity and the unit. On a finite interval
all arguments on the right are at most the point being evaluated.
For group-valued $f$, expand integer multiples by repeated addition
and inverse and use the same finite reindexing.
\end{proof}
In particular $\tau=\mathbf1*\mathbf1$ counts the divisors, and
$\mu*\tau=\mathbf1$. Also $\sum_{d\mid n}\varphi(d)=n$. The integers
$1\le k\le n$ with $n/\gcd(k,n)=d$ correspond to the residues coprime
to $d$ by dividing $k$ by $n/d$. Inversion gives
$\varphi(n)=\sum_{d\mid n}\mu(d)n/d
 =n\prod_{p\mid n}(1-1/p)$, by expansion over the prime masks.
Clearing denominators gives the integer formulation.

For $N\ge1$, group the pairs $dk\le N$ by their product, or by $d$, to get
\[
 1=\sum_{d\le N}\mu(d)\lfloor N/d\rfloor
  =\sum_{k\le N}M(\lfloor N/k\rfloor),\quad
 M(0)=0,\quad M(N)=1-\sum_{k=2}^N M(\lfloor N/k\rfloor).
\]
The prime masks in $\sum_{d^2\mid n}\mu(d)$ give one factor $1-1$ for
every $p^2\mid n$ and an empty product otherwise. Consequently
\[
  S(n)=\sum_{d^2\mid n}\mu(d),\qquad
  Q(N):=\sum_{n\le N}S(n)
     =\sum_{d\le\lfloor\sqrt N\rfloor}\mu(d)\lfloor N/d^2\rfloor.
\]
For $D=\lfloor\sqrt N\rfloor\ge1$ the floor errors total at most $D$,
and $\sum_{d>D}d^{-2}\le1/D$ by comparison with
$1/(d(d-1))$. Thus $|Q(N)/N-1/\zeta(2)|\le D/N+1/D\le2/D$,
using the Dirichlet identity proved below. This gives squarefree density
$6/\pi^2$ with a modulus. The counts with signs $+$ and $-$ are
$(Q(N)\pm M(N))/2$. Their common density is $3/\pi^2$ by the mean
estimate below.

\paragraph{Occupation-number interpretation.}
A vector of prime occupation numbers with finite support maps to
$n=\prod_p p^{a_p}$ and has energy $\sum a_p\log p=\log n$.
The inverse is the valuation vector. Restricting $a_p$ to $0,1$ selects
the squarefree integers. Their parity sign is $(-1)^{\sum a_p}$.
Assigning zero to excluded integers gives exactly $\mu$.
For a cutoff $n\le N$ this is an explicit bijection of finite tables.
Removing the cutoff under the absolutely convergent weights $n^{-s}$
gives the partition sums $\zeta(s)$, $\zeta(s)/\zeta(2s)$, and
$1/\zeta(s)$ in the unrestricted, restricted, and signed cases.
